\chapter{Технологический раздел}

Технологический раздел посвящен практической реализации гибридного метода восстановления глубины из монокулярных изображений.
В данном разделе рассматриваются ключевые аспекты программной реализации и тестирования разработанного решения.

Особое внимание уделяется обоснованию выбора программных средств и библиотек, необходимых для эффективной реализации алгоритмов.

Для объективной оценки качества реконструкции будут формализованы метрики: MAE, RMSE, REL и точность δ, позволяющие количественно оценивать эффективность метода на разных типах изображений.

\section{Выбор программных средств}

Для реализации методов восстановления глубины из монокулярных изображений был выбран язык программирования Python.
Этот выбор обусловлен наличием в его экосистеме всех необходимых библиотек и инструментов для комплексного решения поставленной задачи.

% TODO добавить ссылки на источники

\begin{enumerate}
    \item \textbf{Библиотеки обработки изображений:} OpenCV~\cite{opencv_library}, scikit-image~\cite{van2014scikit} и PIL/Pillow~\cite{clark2015pillow} — обеспечивают функционал для эффективной предобработки архивных фотографий, включая фильтрацию, сегментацию и работу с различными форматами файлов.
    \item \textbf{Библиотеки машинного обучения:} PyTorch~\cite{paszke2019pytorch}, TensorFlow/Keras~\cite{abadi2016tensorflow} и scikit-learn~\cite{pedregosa2011scikit} — предоставляют инструменты для применения предобученной модели MiDaS для оценки глубины, что исключает необходимость обучения собственной нейронной сети и повышает эффективность разработки.
    \item \textbf{Библиотеки научных вычислений:} NumPy~\cite{harris2020array}, SciPy~\cite{virtanen2020scipy}, Matplotlib~\cite{hunter2007matplotlib} и Open3D~\cite{zhou2018open3d} — обеспечивают эффективную работу с многомерными массивами, методы оптимизации, визуализацию результатов и обработку трехмерных данных.
\end{enumerate}

Таким образом, совокупность доступных в Python библиотек делает его достаточным инструментом для решения всех аспектов задачи восстановления глубины, включая предобработку данных, использование готовой модели MiDaS и визуализацию результатов.

\section{Описание входных и выходных данных}

Входными данными для алгоритмов восстановления глубины являются монокулярные цветные фотографии в различных форматах (JPEG, PNG, TIFF).

Результатом работы алгоритмов восстановления глубины являются карта глубин в виде одноканальных изображений, где:

\begin{itemize}
    \item значение пикселя соответствует относительному расстоянию от камеры (меньшие значения — ближе к камере);
    \item диапазон значений нормализован от 0 до 1 (или от 0 до 255 для 8-битных изображений);
    \item разрешение карты глубины соответствует разрешению входного изображения.
\end{itemize}


